\section{Rows, cyclic modules, and derivations}
\label{sec:preliminaries}

This section supplies two inputs to the proof: a uniform implementer
for a rectangular derivation with cyclic domain, and the per-algebra
derivation criterion that will conclude similarity. The cyclic estimate
will first be used to extend control from a small free-product corner
to an entire right-supported subspace.

All Hilbert spaces and algebras are complex, and inner products are linear
in the first variable. All algebra homomorphisms are complex-linear, and
all \(*\)-representations are unital unless stated otherwise. No
separability or faithfulness is implicit. We write \(E_1\) for the closed
unit ball of a normed space \(E\). For a unital self-adjoint operator
algebra \(P\subset\mathcal B(K)\) and \(Y\in\mathcal B(K)\), set
\begin{equation}
 g_P(Y)=\sup_{a\in P_1}\|Ya-aY\|.
 \label{eq:commutator-seminorm}
\end{equation}
The same notation will be used when \(P\) is a specified represented
copy of an abstract algebra.

For a linear map \(u:A\to\mathcal B(K,F)\), its amplification
\(u_{p,q}\) acts entrywise on \(M_{p,q}(A)\), with values in
\(\mathcal B(K^q,F^p)\). Rectangular matrix norms on \(A\) are the usual
\(C^*\)-matrix norms. We abbreviate \(u_{h,h}\) to \(u_h\), and set
\[
 \|u\|_{\mathrm{cb}}=\sup_{h\ge1}\|u_h\|,\qquad
 \|u\|_{\mathrm{row}}=\sup_{h\ge1}\|u_{1,h}\|,\qquad
 \|u\|_{\mathrm{col}}=\sup_{h\ge1}\|u_{h,1}\|.
\]
A map is completely bounded when its cb norm is finite. For an operator
\(Y\), the notation \(Y^{(h)}\) means the diagonal operator on \(h\)
copies of its Hilbert space.

If \((P,\tau)\) is a finite von Neumann algebra with a faithful normal
tracial state, \(L^2(P)=L^2(P,\tau)\) is the completion of \(P\) for
\(\langle x,y\rangle=\tau(y^*x)\). We write \(\|\cdot\|_2\) for this
norm, and \(\|x\|_1=\tau(|x|)\) for \(x\in P\). Left and right
multiplication are denoted by \(L_a\) and \(R_a\), with
\(R_a\xi=\xi a\); thus \(R_aR_b=R_{ba}\). The tracial commutation
theorem identifies \(L(P)'\) with \(\{R_a:a\in P\}\). The support
notation for an \(L^2\)-vector is understood through its affiliated
operator: if \(\xi=v|\xi|\) is its polar decomposition, its left and
right support projections are \(vv^*\) and \(v^*v\), respectively.
They are equivalent in \(P\), and in particular have the same trace.
For a projection \(r\in P\), the condition that \(R_b\) vanish on
\(L^2(P)(1-r)\) is exactly \(b=rb\).
We use these standard facts, and standard-form theory for von Neumann
algebras, without separability restrictions
\citep{Takesaki1979,Haagerup1975}.

\subsection{Uniform row and column estimates}

Dickson's argument derives the row estimate from Haagerup's little
Grothendieck inequality and credits Christensen with the estimate in
this form \citep[Lemmas~2.8--2.9]{Dickson2014}; this is part of the
factorization approach to similarity developed by Pisier
\citep{Pisier1999Degree,Pisier2001Remarks}.
The estimate states that every bounded unital homomorphism
\(\varphi:A\to\mathcal B(K)\), for any unital \(C^*\)-algebra \(A\),
satisfies
\begin{equation}
 \|\varphi\|_{\mathrm{row}}
       \le\sqrt2\,\|\varphi\|^2
 \qquad\text{\citep[Lemma~2.9, Equation~(19)]{Dickson2014}}.
 \label{eq:dickson-row}
\end{equation}
The map \(\varphi^\sharp(a)=\varphi(a^*)^*\) is complex-linear and has
the same norm as \(\varphi\). It is again a homomorphism, since
\[
 \varphi^\sharp(ab)
 =\bigl(\varphi(b^*)\varphi(a^*)\bigr)^*
 =\varphi^\sharp(a)\varphi^\sharp(b).
\]
Apply \eqref{eq:dickson-row} to \(\varphi^\sharp\) on the row
\([a_1^*,\ldots,a_h^*]\), and take the rectangular adjoint. This gives
the same bound for \(\|\varphi\|_{\mathrm{col}}\). Consequently we fix
the absolute constant
\begin{equation}
 c_{\rm row}=4\sqrt2,
 \label{eq:row-constant}
\end{equation}
which bounds both norms whenever \(\|\varphi\|\le2\).

Let \(A\) be a unital \(C^*\)-algebra. A \emph{rectangular derivation}
for \(*\)-representations
\(\sigma:A\to\mathcal B(K)\) and
\(\lambda:A\to\mathcal B(F)\) is a linear map
\(\Delta:A\to\mathcal B(K,F)\) satisfying
\begin{equation}
 \Delta(ab)=\Delta(a)\sigma(b)+\lambda(a)\Delta(b)
 \qquad(a,b\in A).
 \label{eq:rectangular-derivation}
\end{equation}
The derivation identity and unitality imply \(\Delta(1)=0\).

\begin{lemma}[Row and column bounds]
\label{lem:row-column}
Every bounded rectangular derivation satisfies
\begin{equation}
 \max\{\|\Delta\|_{\mathrm{row}},\|\Delta\|_{\mathrm{col}}\}
       \le c_{\rm row}\|\Delta\|.
 \label{eq:rectangular-row-column}
\end{equation}
\end{lemma}

\begin{proof}
The assertion is immediate when \(\Delta=0\). Otherwise put
\(d=\|\Delta\|>0\), and consider the upper triangular map
\begin{equation}
 \Phi(a)=
 \begin{pmatrix}
  \lambda(a)&d^{-1}\Delta(a)\\
  0&\sigma(a)
 \end{pmatrix}
 \in\mathcal B(F\oplus K).
 \label{eq:triangular-homomorphism}
\end{equation}
It is unital and multiplicative by \eqref{eq:rectangular-derivation}.
Its diagonal map is contractive, and its off-diagonal map has norm one;
hence \(\|\Phi\|\le2\). The off-diagonal rectangular compression of
\(\Phi_{p,q}\), from \(K^q\) to \(F^p\), is
\(d^{-1}\Delta_{p,q}\). The row and column estimates for \(\Phi\)
give \eqref{eq:rectangular-row-column}.
\end{proof}

\subsection{Completely bounded implementation}

We use the quantitative completely bounded similarity theorem of
Paulsen: if a unital homomorphism \(\varphi:A\to\mathcal B(H)\) is
completely bounded, there is an invertible \(S\in\mathcal B(H)\) such
that \(S\varphi(\cdot)S^{-1}\) is completely contractive and
\(\|S\|\|S^{-1}\|\le\|\varphi\|_{\mathrm{cb}}\)
\citep[p.~226]{Paulsen1984}. For a unital \(C^*\)-algebra the
resulting homomorphism is a \(*\)-homomorphism: a unitary and its inverse
are sent to contractions, so the image is a unitary; unitaries linearly
span the algebra. For the later cyclic-domain estimate, set
\begin{equation}
 c_0=(1+4c_{\rm row})^2.
 \label{eq:cyclic-constant}
\end{equation}

\begin{lemma}[Completely bounded rectangular derivations]
\label{lem:cb-implementation}
Let \(\sigma,\lambda,\Delta\) satisfy
\eqref{eq:rectangular-derivation}. If \(\Delta\) is completely bounded,
then some \(V\in\mathcal B(K,F)\) satisfies
\begin{equation}
 \Delta(a)=V\sigma(a)-\lambda(a)V\qquad(a\in A).
 \label{eq:rectangular-implementation}
\end{equation}
If \(d=\|\Delta\|>0\), one may choose
\[
 \|V\|\le d\bigl(1+\|\Delta\|_{\mathrm{cb}}/d\bigr)^2.
\]
\end{lemma}

\begin{proof}
We use the triangular-metric argument recalled in
\citet[Lemma~1.2]{Ozawa2006}, allowing different domain and codomain
representations and keeping track of the implementer norm.
For \(\Delta=0\), take \(V=0\). Otherwise the map \(\Phi\) in
\eqref{eq:triangular-homomorphism} is completely bounded, with
\[
 \|\Phi\|_{\mathrm{cb}}
 \le 1+\|\Delta\|_{\mathrm{cb}}/d=:G.
\]
Choose the similarity \(S\) supplied by Paulsen's theorem, and put
\(D=S^*S\). The fact that \(S\Phi(\cdot)S^{-1}\) preserves adjoints
is equivalent to
\begin{equation}
 D\Phi(a)=\Phi(a^*)^*D\qquad(a\in A).
 \label{eq:intertwining-metric}
\end{equation}
Write \(D=(D_{ij})_{i,j=1}^2\) relative to \(F\oplus K\). Its first
row in \eqref{eq:intertwining-metric} gives
\begin{align}
 D_{11}\lambda(a)&=\lambda(a)D_{11},\notag\\
 d^{-1}D_{11}\Delta(a)+D_{12}\sigma(a)&=\lambda(a)D_{12}.
 \label{eq:metric-blocks}
\end{align}
Since \(D\ge\|D^{-1}\|^{-1}I\), its compression \(D_{11}\) is
positive invertible on \(F\), with
\(\|D_{11}^{-1}\|\le\|D^{-1}\|\). Its inverse also commutes with
\(\lambda(A)\). Hence
\[
 V=-dD_{11}^{-1}D_{12}
\]
satisfies \eqref{eq:rectangular-implementation}, by
\eqref{eq:metric-blocks}, and
\[
 \|V\|
 \le d\,\|D^{-1}\|\|D\|
 =d\,\bigl(\|S^{-1}\|\|S\|\bigr)^2
 \le dG^2.
\]
If either Hilbert space is zero, \(\Delta=0\), so it has already been
covered.
\end{proof}

\subsection{Cyclic domains}

The cyclic hypothesis needed below concerns only the domain
\(*\)-representation. The output representation may have arbitrary
multiplicity.

\begin{lemma}[Column factorization in a cyclic representation]
\label{lem:cyclic-factorization}
Let \(P\subset\mathcal B(K)\) be a unital von Neumann algebra with a
cyclic vector. For every integer \(h\ge1\) and every unit vector
\(\xi=(\xi_1,\ldots,\xi_h)\in K^h\), there are a unit vector
\(\zeta\in K\) and a contraction \(v=(v_i)\in M_{h,1}(P)\) such
that \(\xi_i=v_i\zeta\) for every \(i\).
\end{lemma}

\begin{proof}
Put \(N=P'\), and choose a cyclic vector \(\Omega\) for \(P\).
It separates \(N\): if \(x\in N\) and \(x\Omega=0\), then
\(xa\Omega=ax\Omega=0\) for every \(a\in P\), whence \(x=0\).
Let \(L=\overline{N\Omega}\subset K\). This subspace reduces \(N\),
and restriction \(r:N\to\mathcal B(L)\) is faithful and normal.
Indeed, an element in its kernel annihilates \(\Omega\); normality
follows from compression to a reducing subspace. Thus it is a normal
\(*\)-isomorphism onto the von Neumann algebra \(r(N)\), and its
inverse is normal. The vector \(\Omega\) is cyclic and separating for
\(r(N)\).

The standard-form theorem for a representation with a cyclic separating
vector realizes every normal positive functional on \(r(N)\) as a
vector functional, with its vector in the natural cone
\citep[Lemma~2.10(1), pp.~278--279]{Haagerup1975}.
Apply it, through the normal isomorphism \(r\),
to
\[
 \omega(x)=\sum_{i=1}^h\langle x\xi_i,\xi_i\rangle,
 \qquad x\in N.
\]
We obtain \(\zeta\in L\subset K\) such that
\(\langle x\zeta,\zeta\rangle=\omega(x)\) for every \(x\in N\).
In particular \(\|\zeta\|^2=\omega(1)=1\). The assignment
\[
 x\zeta\longmapsto(x\xi_1,\ldots,x\xi_h),\qquad x\in N,
\]
is isometric, since both squared norms equal \(\omega(x^*x)\).
It extends to an isometry from \(\overline{N\zeta}\) into \(K^h\).
Extend it by zero on \(\overline{N\zeta}^{\perp}\), obtaining a
contraction \(v:K\to K^h\). The subspace \(\overline{N\zeta}\)
reduces \(N\), and the isometry intertwines the \(N\)-actions.
Therefore \(vx=x^{(h)}v\) for every \(x\in N\). Each coordinate
\(v_i\) lies in \(N'=P\), and \(v\zeta=\xi\), as required.
All these constructions, including the standard-form assertion, are
valid on nonseparable Hilbert spaces.
\end{proof}

\begin{lemma}[Implementation on a cyclic domain]
\label{lem:cyclic-implementation}
Let \(A\) be a unital \(C^*\)-algebra, let
\(\sigma:A\to\mathcal B(K)\) and
\(\lambda:A\to\mathcal B(F)\) be \(*\)-representations, and suppose
that \(\sigma\) has a cyclic vector. Every bounded rectangular
derivation \(\Delta:A\to\mathcal B(K,F)\) is completely bounded,
with
\begin{equation}
 \|\Delta\|_{\mathrm{cb}}\le4c_{\rm row}\|\Delta\|.
 \label{eq:cyclic-cb}
\end{equation}
Moreover there is \(V\in\mathcal B(K,F)\) such that
\begin{equation}
 \Delta(a)=V\sigma(a)-\lambda(a)V\quad(a\in A),\qquad
 \|V\|\le c_0\|\Delta\|.
 \label{eq:cyclic-implementation}
\end{equation}
\end{lemma}

\begin{proof}
Zero spaces and the zero derivation are immediate, so assume they do
not occur. Put \(P=\sigma(A)''\). Fix \(h\ge1\),
\(X\in M_h(A)_1\), and a unit vector \(\xi\in K^h\).
Lemma~\ref{lem:cyclic-factorization} gives a unit vector \(\zeta\in K\)
and a contraction \(v\in M_{h,1}(P)\) with \(v\zeta=\xi\).

Matrix Kaplansky density gives a net of column contractions
\(z_\alpha\in M_{h,1}(\sigma(A))\) converging strongly to \(v\).
To see the rectangular assertion directly, regard \(v\) as the first
column of a matrix in \(M_h(P)\), with all other columns zero.
Approximate that matrix strongly by contractions in
\(M_h(\sigma(A))\), and compress to its first input coordinate.
The \(*\)-image \(\sigma(A)\) is the \(C^*\)-quotient
\(A/\ker\sigma\), and the same quotient identification holds at each
square matrix level. Consequently each \(z_\alpha\) lifts to
\(w_\alpha\in M_{h,1}(A)\) with
\begin{equation}
 \|w_\alpha\|\le2,\qquad
 \sigma_{h,1}(w_\alpha)=z_\alpha,\qquad
 \sigma_{h,1}(w_\alpha)\zeta\longrightarrow\xi.
 \label{eq:cyclic-column-lifts}
\end{equation}
For example, first lift the square matrix whose first column is
\(z_\alpha\) and whose other columns vanish to a matrix of norm less
than \(2\), then compress its first column. Only fixed slack in the
quotient norm is being used.

Entrywise application of \eqref{eq:rectangular-derivation} gives
\[
 \Delta_h(X)\sigma_{h,1}(w_\alpha)
   =\Delta_{h,1}(Xw_\alpha)
       -\lambda_h(X)\Delta_{h,1}(w_\alpha).
\]
Since \(\lambda\) is completely contractive, Lemma~\ref{lem:row-column}
and \eqref{eq:cyclic-column-lifts} imply
\[
 \bigl\|\Delta_h(X)\sigma_{h,1}(w_\alpha)\zeta\bigr\|
 \le c_{\rm row}\|\Delta\|
       \bigl(\|Xw_\alpha\|+\|\lambda_h(X)\|\|w_\alpha\|\bigr)
 \le4c_{\rm row}\|\Delta\|.
\]
At the fixed finite level \(h\), \(\Delta_h(X)\) is already a bounded
operator. We may therefore take the limit in
\eqref{eq:cyclic-column-lifts}, obtaining the same bound on
\(\|\Delta_h(X)\xi\|\). Taking the supremum over \(\xi,X,h\)
proves \eqref{eq:cyclic-cb}. The quotient argument did not require
\(\Delta\) to vanish on \(\ker\sigma\); its identity was applied to
the chosen lifts in \(A\). Lemma~\ref{lem:cb-implementation} and
\eqref{eq:cyclic-constant} now give
\eqref{eq:cyclic-implementation}.
\end{proof}

\subsection{The derivation criterion for similarity}

For the last step of the proof we use Kirchberg's equivalence in its
per-algebra form. We state the quantifiers because the final
\(*\)-representation used to test derivations is independent of the
bounded homomorphism in the similarity problem.

\begin{theorem}[Kirchberg's similarity--derivation equivalence]
\label{thm:kirchberg}
For a unital \(C^*\)-algebra \(A\), the following properties are
equivalent.
\begin{enumerate}
 \item Every bounded unital homomorphism
 \(\pi:A\to\mathcal B(H)\), on any Hilbert space \(H\), is similar
 by a bounded invertible operator to a \(*\)-homomorphism.
 \item For every faithful unital \(*\)-representation
 \(\sigma:A\to\mathcal B(K)\), on any Hilbert space \(K\), every
 bounded linear map \(\Delta:A\to\mathcal B(K)\) satisfying
 \[
  \Delta(ab)=\Delta(a)\sigma(b)+\sigma(a)\Delta(b)
 \]
 has the form \(\Delta(a)=V\sigma(a)-\sigma(a)V\) for some
 \(V\in\mathcal B(K)\).
\end{enumerate}
\end{theorem}

This is the theorem of \citet{Kirchberg1996}; the exact formulation
with every faithful representation is stated in
\citet[Theorem~1.1]{Ozawa2006}. The usual derivation property is stated
without a continuity hypothesis. Ringrose's automatic continuity
theorem for derivations from \(C^*\)-algebras into Banach bimodules
\citep[Theorem~2, p.~434]{Ringrose1972} makes it equivalent to the
bounded version above.
The triangular metric computation in
Lemma~\ref{lem:cb-implementation} uses only Paulsen's established
completely bounded theorem; the unrestricted similarity conclusion
will follow from Theorem~\ref{thm:kirchberg} after all bounded
derivations have been treated.
